\clearpage
\chapter{Marco normativo, ética y protección de datos}

Se analiza el entorno legal y ético aplicable, con especial énfasis en el Reglamento General de Protección de Datos (RGPD) para el uso de datos masivos y las directrices éticas sobre el uso de la Inteligencia Artificial en el sector financiero.

\section{Identificación del tratamiento de datos}

En el desarrollo del presente Trabajo de Fin de Máster, el tratamiento de datos se limita estrictamente a series temporales de precios de activos financieros públicos procedentes del índice S\&P 500. Se hace constar que no se han utilizado, procesado ni almacenado datos de carácter personal de terceros identificados o identificables, cumpliendo con el principio de minimización de datos.

\section{Cumplimiento del RGPD y la Ley Orgánica 3/2018 (LOPDGDD)}

A pesar de la naturaleza pública de los datos financieros, el proyecto se alinea con el Reglamento General de Protección de Datos (RGPD) y la Ley Orgánica 3/2018 en los siguientes términos:

\begin{itemize}
    \item \textbf{Finalidad:} Los datos son tratados con el fin exclusivo de investigación académica y desarrollo de modelos de aprendizaje profundo.
    \item \textbf{Seguridad:} El procesamiento de los datos se realiza en un entorno local, garantizando la integridad de los mismos y evitando cualquier fuga de información que pudiera afectar a la infraestructura del sistema.
    \item \textbf{Análisis de Riesgo:} Dado que no existe tratamiento de datos de personas físicas, el riesgo para los derechos y libertades de las personas es nulo.
\end{itemize}

\section{Consideraciones éticas y propiedad intelectual}

El acceso a los datos se ha realizado mediante la API de Yahoo Finance, respetando los términos y condiciones de uso para fines educativos y de investigación. El código fuente desarrollado y los escenarios sintéticos generados por la arquitectura VAE-GAN son de autoría propia, garantizando la transparencia en el proceso de generación de información. Sin embargo, como suele ocurrir con este tipo de modelos, la elección de los activos puede carecer de explicabilidad aparente.

Se hace constar que la elección de unos activos determinados frente a otros puede no ser explicable mediante lógicas económicas o financieras fundamentales, basadas en las métricas principales en este caso de las empresas seleccionadas. Se recomienda la comprensión profunda de la naturaleza del mercado.

Incluso con la capacidad de construir un modelo sólido, no se puede aseverar que la cartera vaya realmente a cumplir con las expectativas de rentabilidad calculadas. 